% 由 tools/md2tex.py 自 2019_数学真题_文_试题.md 生成，勿手改（改源 MD 后重跑）


\begin{gknotice}
\textbf{ 卷种提示}：本卷为文理分科时代的旧结构，含\textbf{现已取消}的选做题（第 22/23 题）与\textbf{已删除}的程序框图考点（第 9 题）。适合\textbf{按知识点抽题}，不适合整套限时模拟。

\textbf{与同年理科卷重合}：第 3、4、5、8、9、12、13 题及第 19(1)、22、23 题与理科卷完全相同。
\end{gknotice}

\par\medskip

\gkpart{一、选择题}

\begin{gktip}
本题共 12 小题，每小题 5 分，共 60 分。
\end{gktip}

\gkti{1}设 $z=\dfrac{3-\mathrm i}{1+2\mathrm i}$，则 $|z|=$（　　）

\begin{gkopts}(4)
\task A．$2$
\task B．$\sqrt3$
\task C．$\sqrt2$
\task D．$1$
\end{gkopts}

\gkti{2}已知集合 $U=\{1,2,3,4,5,6,7\}$，$A=\{2,3,4,5\}$，$B=\{2,3,6,7\}$，则 $B\cap\complement_U A=$（　　）

\begin{gkopts}(4)
\task A．$\{1,6\}$
\task B．$\{1,7\}$
\task C．$\{6,7\}$
\task D．$\{1,6,7\}$
\end{gkopts}

\gkti{3}已知 $a=\log_2 0.2$，$b=2^{0.2}$，$c=0.2^{0.3}$，则（　　）

\begin{gkopts}(4)
\task A．$a<b<c$
\task B．$a<c<b$
\task C．$c<a<b$
\task D．$b<c<a$
\end{gkopts}

\gkti{4}古希腊时期，人们认为最美人体的头顶至肚脐的长度与肚脐至足底的长度之比是 $\dfrac{\sqrt5-1}2\ (\approx0.618$，称为黄金分割比例$)$，著名的"断臂维纳斯"便是如此。此外，最美人体的头顶至咽喉的长度与咽喉至肚脐的长度之比也是 $\dfrac{\sqrt5-1}2$。若某人满足上述两个黄金分割比例，且腿长为 $105\ \mathrm{cm}$，头顶至脖子下端的长度为 $26\ \mathrm{cm}$，则其身高可能是（　　）

\begin{gkopts}(2)
\task A．$165\ \mathrm{cm}$
\task B．$175\ \mathrm{cm}$
\task C．$185\ \mathrm{cm}$
\task D．$190\ \mathrm{cm}$
\end{gkopts}

\gkti{5}函数 $f(x)=\dfrac{\sin x+x}{\cos x+x^{2}}$ 在 $[-\pi,\pi]$ 的图象大致为（　　）

\begin{center}% 2019 全国Ⅰ卷 理5/文5：f(x)=(sin x+x)/(cos x+x^2) 图象四选一
% (C)(D) 由精确表达式绘制（D=f 为正确图象，C=-f）；
% (A)(B) 为干扰项，依据原卷形态用样条示意重绘（原卷亦为定性示意图）
\begin{tikzpicture}[font=\footnotesize]
\pgfplotsset{
  panel/.style={width=4.6cm, height=3.4cm, axis lines=middle,
    xmin=-3.6, xmax=3.6, ymin=-1.6, ymax=1.7,
    xtick={-3.14159,3.14159}, xticklabels={$-\pi$,$\pi$},
    ytick={1}, yticklabels={$1$},
    xlabel={$x$}, ylabel={$y$},
    every axis x label/.style={at={(current axis.right of origin)}, anchor=west},
    every axis y label/.style={at={(current axis.above origin)}, anchor=south}}}
  \begin{axis}[panel, at={(0,0)}, name=axA]
    \addplot[smooth, thick] coordinates
      {(-3.14,0.55) (-2.2,0.42) (-1.2,0.18) (-0.5,-0.15) (0,0)
       (0.45,0.75) (0.8,1.05) (1.6,0.8) (2.4,0.4) (3.14,0.1)};
    \node[below left] at (axis cs:0,0) {$O$};
  \end{axis}
  \node[below=0.5mm] at (axA.south) {(A)};
  \begin{axis}[panel, at={(5.5cm,0)}, name=axB]
    \addplot[smooth, thick] coordinates
      {(-3.14,-0.02) (-2.4,-0.25) (-1.5,-0.45) (-0.6,-0.3) (0,0)
       (0.7,0.45) (1.3,0.58) (2.2,0.35) (3.14,0.12)};
    \node[below left] at (axis cs:0,0) {$O$};
  \end{axis}
  \node[below=0.5mm] at (axB.south) {(B)};
  \begin{axis}[panel, at={(0,-4.4cm)}, name=axC]
    \addplot[domain=-3.14159:3.14159, samples=200, thick]
      {-(sin(deg(x))+x)/(cos(deg(x))+x^2)};
    \node[below left] at (axis cs:0,0) {$O$};
  \end{axis}
  \node[below=0.5mm] at (axC.south) {(C)};
  \begin{axis}[panel, at={(5.5cm,-4.4cm)}, name=axD]
    \addplot[domain=-3.14159:3.14159, samples=200, thick]
      {(sin(deg(x))+x)/(cos(deg(x))+x^2)};
    \node[below left] at (axis cs:0,0) {$O$};
  \end{axis}
  \node[below=0.5mm] at (axD.south) {(D)};
\end{tikzpicture}
\end{center}

\gkti{6}某学校为了解 $1000$ 名新生的身体素质，将这些学生编号为 $1,2,\cdots,1000$，从这些新生中用系统抽样方法等距抽取 $100$ 名学生进行体质测验。若 $46$ 号学生被抽到，则下面 $4$ 名学生中被抽到的是（　　）

\begin{gkopts}(4)
\task A．$8$ 号学生
\task B．$200$ 号学生
\task C．$616$ 号学生
\task D．$815$ 号学生
\end{gkopts}

\gkti{7}$\tan255^\circ=$（　　）

\begin{gkopts}(4)
\task A．$-2-\sqrt3$
\task B．$-2+\sqrt3$
\task C．$2-\sqrt3$
\task D．$2+\sqrt3$
\end{gkopts}

\gkti{8}已知非零向量 $\boldsymbol a$，$\boldsymbol b$ 满足 $|\boldsymbol a|=2|\boldsymbol b|$，且 $(\boldsymbol a-\boldsymbol b)\perp\boldsymbol b$，则 $\boldsymbol a$ 与 $\boldsymbol b$ 的夹角为（　　）

\begin{gkopts}(4)
\task A．$\dfrac\pi6$
\task B．$\dfrac\pi3$
\task C．$\dfrac{2\pi}3$
\task D．$\dfrac{5\pi}6$
\end{gkopts}

\gkti{9}如图是求 $\dfrac1{2+\dfrac1{2+\dfrac12}}$ 的程序框图，图中空白框中应填入（　　）

\begin{center}% 2019 全国Ⅰ卷 理8/文9：求连分数 1/(2+1/(2+1/2)) 的程序框图
% 空白矩形为题目要求填入的位置（正确答案 A=1/(2+A)）
\begin{tikzpicture}[font=\small, >=Stealth, node distance=7mm,
  box/.style={draw, minimum height=1.9em, inner xsep=7pt},
  term/.style={draw, rounded corners=7pt, minimum height=1.9em, inner xsep=9pt},
  io/.style={draw, trapezium, trapezium left angle=70, trapezium right angle=110,
             minimum height=1.9em, inner xsep=2pt},
  dec/.style={draw, diamond, aspect=2.4, inner sep=1.5pt}]
  \node[term] (start) {开始};
  \node[box, below=of start] (init) {$A=\dfrac12$};
  \node[box, below=of init] (k1) {$k=1$};
  \coordinate[below=4mm of k1] (merge);
  \node[dec, below=2mm of merge] (cond) {$k\le2$};
  \node[box, below=of cond, minimum width=5.6em] (blank) {\phantom{$A=\dfrac1{2+A}$}};
  \node[box, below=of blank] (inc) {$k=k+1$};
  \node[io, right=14mm of cond] (out) {输出 $A$};
  \node[term, below=of out] (end) {结束};
  \draw[->] (start) -- (init);
  \draw[->] (init) -- (k1);
  \draw[->] (k1) -- (cond);
  \draw[->] (cond) -- node[left]{是} (blank);
  \draw[->] (blank) -- (inc);
  \draw[->] (inc.west) -- ++(-8mm,0) |- (merge);
  \draw[->] (cond) -- node[above]{否} (out);
  \draw[->] (out) -- (end);
\end{tikzpicture}
\end{center}

\begin{gkopts}(2)
\task A．$A=\dfrac1{2+A}$
\task B．$A=2+\dfrac1A$
\task C．$A=\dfrac1{1+2A}$
\task D．$A=1+\dfrac1{2A}$
\end{gkopts}

\begin{gktip}
\textbf{注}：程序框图（算法初步）在现行新课标中已删除。
\end{gktip}

\gkti{10}双曲线 $C:\dfrac{x^{2}}{a^{2}}-\dfrac{y^{2}}{b^{2}}=1\ (a>0,\ b>0)$ 的一条渐近线的倾斜角为 $130^\circ$，则 $C$ 的离心率为（　　）

\begin{gkopts}(2)
\task A．$2\sin40^\circ$
\task B．$2\cos40^\circ$
\task C．$\dfrac1{\sin50^\circ}$
\task D．$\dfrac1{\cos50^\circ}$
\end{gkopts}

\gkti{11}$\triangle ABC$ 的内角 $A$，$B$，$C$ 的对边分别为 $a$，$b$，$c$，已知 $a\sin A-b\sin B=4c\sin C$，$\cos A=-\dfrac14$，则 $\dfrac bc=$（　　）

\begin{gkopts}(4)
\task A．$6$
\task B．$5$
\task C．$4$
\task D．$3$
\end{gkopts}

\gkti{12}已知椭圆 $C$ 的焦点为 $F_1(-1,0)$，$F_2(1,0)$，过 $F_2$ 的直线与 $C$ 交于 $A$，$B$ 两点。若 $\left|AF_2\right|=2\left|F_2B\right|$，$|AB|=\left|BF_1\right|$，则 $C$ 的方程为（　　）

\begin{gkopts}(2)
\task A．$\dfrac{x^{2}}2+y^{2}=1$
\task B．$\dfrac{x^{2}}3+\dfrac{y^{2}}2=1$
\task C．$\dfrac{x^{2}}4+\dfrac{y^{2}}3=1$
\task D．$\dfrac{x^{2}}5+\dfrac{y^{2}}4=1$
\end{gkopts}

\par\medskip

\gkpart{二、填空题}

\begin{gktip}
本题共 4 小题，每小题 5 分，共 20 分。
\end{gktip}

\gkti{13}曲线 $y=3\left(x^{2}+x\right)\mathrm e^{x}$ 在点 $(0,0)$ 处的切线方程为 \gkblank{}．

\gkti{14}记 $S_n$ 为等比数列 $\left\{a_n\right\}$ 的前 $n$ 项和。若 $a_1=1$，$S_3=\dfrac34$，则 $S_4=$ \gkblank{}．

\gkti{15}函数 $f(x)=\sin\left(2x+\dfrac{3\pi}2\right)-3\cos x$ 的最小值为 \gkblank{}．

\gkti{16}已知 $\angle ACB=90^\circ$，$P$ 为平面 $ABC$ 外一点，$PC=2$，点 $P$ 到 $\angle ACB$ 两边 $AC$，$BC$ 的距离均为 $\sqrt3$，那么 $P$ 到平面 $ABC$ 的距离为 \gkblank{}．

\par\medskip

\gkpart{三、解答题}

\begin{gktip}
共 70 分。第 17～21 题为必考题，第 22、23 题为选考题。
\end{gktip}

\gksubpart{（一）必考题：共 60 分}

\gkti{17}\textbf{（12 分）} 某商场为提高服务质量，随机调查了 $50$ 名男顾客和 $50$ 名女顾客，每位顾客对该商场的服务给出满意或不满意的评价，得到如表列联表：

\begin{center}
\begin{tabular}{ccc}
\toprule
 & 满意 & 不满意 \\
\midrule
男顾客 & $40$ & $10$ \\
女顾客 & $30$ & $20$ \\
\bottomrule
\end{tabular}
\end{center}

  (1) 分别估计男、女顾客对该商场服务满意的概率；

  (2) 能否有 $95\%$ 的把握认为男、女顾客对该商场服务的评价有差异？

  附：$K^{2}=\dfrac{n(ad-bc)^{2}}{(a+b)(c+d)(a+c)(b+d)}$．

\begin{center}
\begin{tabular}{cccc}
\toprule
$P\left(K^{2}\ge k\right)$ & $0.050$ & $0.010$ & $0.001$ \\
\midrule
$k$ & $3.841$ & $6.635$ & $10.828$ \\
\bottomrule
\end{tabular}
\end{center}

\gkti{18}\textbf{（12 分）} 记 $S_n$ 为等差数列 $\left\{a_n\right\}$ 的前 $n$ 项和，已知 $S_9=-a_5$。

  (1) 若 $a_3=4$，求 $\left\{a_n\right\}$ 的通项公式；

  (2) 若 $a_1>0$，求使得 $S_n\ge a_n$ 的 $n$ 的取值范围。

\gkti{19}\textbf{（12 分）} 如图，直四棱柱 $ABCD\text{-}A_1B_1C_1D_1$ 的底面是菱形，$AA_1=4$，$AB=2$，$\angle BAD=60^\circ$，$E$，$M$，$N$ 分别是 $BC$，$BB_1$，$A_1D$ 的中点。

\begin{center}% 2019 全国Ⅰ卷 理18/文19：直四棱柱 ABCD-A1B1C1D1（底面菱形），
% E、M、N 分别为 BC、BB1、A1D 的中点
\begin{tikzpicture}[scale=0.92, font=\small]
  % 底面菱形（斜二测示意）：A 前左、B 前右、C 右后、D 左后
  \coordinate (A) at (0,0);
  \coordinate (B) at (2.1,-0.25);
  \coordinate (C) at (3.15,0.6);
  \coordinate (D) at (1.05,0.85);
  \coordinate (A1) at (0,3.0);
  \coordinate (B1) at (2.1,2.75);
  \coordinate (C1) at (3.15,3.6);
  \coordinate (D1) at (1.05,3.85);
  \coordinate (E) at ($(B)!0.5!(C)$);
  \coordinate (M) at ($(B)!0.5!(B1)$);
  \coordinate (N) at ($(A1)!0.5!(D)$);
  \draw (A) -- (B) -- (C);
  \draw[dashed] (C) -- (D) -- (A);
  \draw (A) -- (A1)  (B) -- (B1)  (C) -- (C1);
  \draw[dashed] (D) -- (D1);
  \draw (A1) -- (B1) -- (C1) -- (D1) -- cycle;
  \draw[dashed] (A1) -- (D) ;
  \draw[dashed] (N) -- (M);
  \draw[dashed] (D) -- (E);
  \draw (C1) -- (E);
  \draw[dashed] (C1) -- (D);
  \node[below left] at (A) {$A$};
  \node[below] at (B) {$B$};
  \node[right] at (C) {$C$};
  \node[below right, inner sep=0.5pt] at (D) {$D$};
  \node[above left, inner sep=1pt] at (A1) {$A_1$};
  \node[below right, inner sep=1pt] at (B1) {$B_1$};
  \node[right] at (C1) {$C_1$};
  \node[above] at (D1) {$D_1$};
  \node[below right, inner sep=1pt] at (E) {$E$};
  \node[right, inner sep=2pt] at (M) {$M$};
  \node[left, inner sep=2.5pt] at (N) {$N$};
\end{tikzpicture}
\end{center}

  (1) 证明：$MN\parallel$ 平面 $C_1DE$；

  (2) 求点 $C$ 到平面 $C_1DE$ 的距离。

\gkti{20}\textbf{（12 分）} 已知函数 $f(x)=2\sin x-x\cos x-x$，$f'(x)$ 为 $f(x)$ 的导数。

  (1) 证明：$f'(x)$ 在区间 $(0,\pi)$ 存在唯一零点；

  (2) 若 $x\in[0,\pi]$ 时，$f(x)\ge ax$，求 $a$ 的取值范围。

\gkti{21}\textbf{（12 分）} 已知点 $A$，$B$ 关于坐标原点 $O$ 对称，$|AB|=4$，$\odot M$ 过点 $A$，$B$ 且与直线 $x+2=0$ 相切。

  (1) 若 $A$ 在直线 $x+y=0$ 上，求 $\odot M$ 的半径；

  (2) 是否存在定点 $P$，使得当 $A$ 运动时，$|MA|-|MP|$ 为定值？并说明理由。

\par\medskip

\gksubpart{（二）选考题：共 10 分，第 22、23 题任选一题作答}

\begin{gktip}
\textbf{ 现行考纲已取消选做题}，以下两题仅作历史资料保留。
\end{gktip}

\gkti{22}\textbf{（10 分）〔选修 4-4：坐标系与参数方程〕} （\textbf{已取消}）

在直角坐标系 $xOy$ 中，曲线 $C$ 的参数方程为 $\begin{cases}x=\dfrac{1-t^{2}}{1+t^{2}}\\[2mm] y=\dfrac{4t}{1+t^{2}}\end{cases}$（$t$ 为参数）。以坐标原点 $O$ 为极点，$x$ 轴的正半轴为极轴建立极坐标系，直线 $l$ 的极坐标方程为 $2\rho\cos\theta+\sqrt3\rho\sin\theta+11=0$。

  (1) 求 $C$ 和 $l$ 的直角坐标方程；

  (2) 求 $C$ 上的点到 $l$ 距离的最小值。

\gkti{23}\textbf{（10 分）〔选修 4-5：不等式选讲〕} （\textbf{已取消}）

已知 $a$，$b$，$c$ 为正数，且满足 $abc=1$。证明：

  (1) $\dfrac1a+\dfrac1b+\dfrac1c\le a^{2}+b^{2}+c^{2}$；

  (2) $(a+b)^{3}+(b+c)^{3}+(c+a)^{3}\ge24$。

\par\medskip

